\subsection{Capon Beamformer (MVDR)}\label{sec:doa:capon}

El beamformer de \textbf{Capon}, también conocido como 
\textbf{MVDR (Minimum Variance Distortionless Response)}, 
se propuso como una mejora directa del método Delay-and-Sum.  
Su idea central es apuntar hacia una dirección deseada sin distorsionar la señal proveniente de allí, 
al mismo tiempo que se minimiza la energía recibida desde todas las demás direcciones.  

De manera intuitiva, si el DAS se asemeja a una linterna con un haz ancho, 
Capon actúa como un láser mucho más fino: concentra la ganancia únicamente en el ángulo de interés 
y suprime la contribución de señales interferentes \cite{YangBook}.

\begin{itemize}
	\item \textbf{Ganancia unitaria en la dirección deseada}: se garantiza que la señal en el ángulo $\varphi$ no se distorsione.  
	\item \textbf{Mínima potencia desde otras direcciones}: los pesos del arreglo se ajustan para cancelar la energía procedente del resto del espacio.  
\end{itemize}

\subsubsection*{Características}
\begin{itemize}
	\item Mayor resolución que DAS: lóbulos principales más estrechos y mejor discriminación angular.  
	\item Rechazo eficaz de interferencias fuera de la dirección de interés.  
	\item Requiere una estimación precisa de la matriz de correlación $R_{uu}$; con pocos datos puede degradarse.  
	\item Más costoso computacionalmente, ya que exige invertir matrices.  
	\item Sensible a errores de calibración en el arreglo.  
\end{itemize}

\subsubsection*{Formulación matemática}

El problema de diseño de Capon se plantea como una optimización:
\begin{equation}
	\min_{\mathbf{w}} \; \mathbf{w}^H R_{uu}\, \mathbf{w}
	\quad \text{sujeto a} \quad \mathbf{w}^H \mathbf{a}(\varphi) = 1,
\end{equation}
donde $R_{uu}$ es la matriz de correlación de las señales, 
$\mathbf{a}(\varphi)$ es el \emph{steering vector}, 
y $\mathbf{w}$ los pesos del beamformer.

La solución, obtenida con multiplicadores de Lagrange, es:
\begin{equation}
	\mathbf{w}_c(\varphi) = \frac{R_{uu}^{-1}\, \mathbf{a}(\varphi)}
	{\mathbf{a}^H(\varphi)\, R_{uu}^{-1}\, \mathbf{a}(\varphi)}.
\end{equation}

El espectro espacial asociado al beamformer de Capon es:
\begin{equation}
	P_{\text{Capon}}(\varphi) = 
	\frac{1}{\mathbf{a}^H(\varphi)\, R_{uu}^{-1}\, \mathbf{a}(\varphi)}.
\end{equation}

Los máximos de $P_{\text{Capon}}(\varphi)$ señalan las direcciones estimadas de arribo, 
con lóbulos mucho más estrechos que los de DAS, 
tal como se aprecia en la Fig.~\ref{fig:doa:caponMatlab}.

\insertarfigura{fig/CAPON.png}{Espectro espacial obtenido con el beamformer de Capon (MVDR). 
	Se simularon tres fuentes incidentes en $30^\circ$, $50^\circ$ y $70^\circ$, 
	con una relación señal-ruido (SNR) de 5~dB, usando un arreglo de 10 sensores 
	con espaciamiento normalizado $d_{norm}=0,5$. 
	En comparación con DAS, los lóbulos principales son más estrechos, 
	lo que permite una mejor discriminación angular entre fuentes cercanas.}{fig:doa:caponMatlab}{1}
